\chapter{为什么网络通信需要持久端点}

\begin{statebox}
\textbf{继承的机器。} 内核已经能把异步设备请求组织成长期对象，也能用 fd 引用文件。
单个网络数据报仍不能表达长期会话：地址由谁占用、连接处于哪个阶段、发送者何时需要背压、
接收者等待什么条件、关闭后迟到的数据属于谁，都需要独立保存。
\end{statebox}

\inputchapterbody{chapters/ch18-socket-network-entry}

\section{网络事件、流量控制与提交边界}
\inputdetail{chapters/ch18-network-persistence-deep-dive}

\begin{problemframe}
\textbf{尚未解决的问题。} 地址空间、fd 与 socket 已经隔离了对象命名，却没有回答某个
主体是否应当获得某项操作权限，也不能限制一组进程耗尽全机资源。下一章从一次主体访问
一个对象的判定开始建立最小权限。
\end{problemframe}
